\newpage
\begin{abstractspacing}

\begin{center}
    \textbf{Abstract}
\end{center}

\noindent\textbf{Short version}

AI has the potential, this project claims, to help humanity become a technologically mature species; in this context, is it meaningful to ask "what values ought to be transmitted to our successors?" The written component (a) reviews the historical context (particularly from expected utility theory) and current concerns of "AI alignment", as well as certain adjacencies ((1) an alleged discursive substructure of Judeo-Christian eschatology; (2) overlaps with, and critiques from, neorationalist philosophy within the CCRU lineage); (b) examines the variety of possible cognitions as a way of decentring anthropocentric biases; (c) analyzes subject-neutral axiology (the "view from nowhere") and its criticisms. Ambitiously, it considers how acausal decision theories might result in conceptions of morality that better suit space-faring AIs.

The practice component involves collaborating with AI alignment researchers and using large language models (LLMs) and reinforcement learners (RLs) as research partners and experimental subjects. This includes:

\begin{enumerate}
\item Analyzing how LLMs approach Newcomb's Paradox to gain insight into how agents with strong mutual predictive abilities might reason about each other's actions 
\item Using LLM reasoning on acausal decision theory to inform how likely advanced AI may be to behave benevolently towards humans
\item Testing if LLMs incorporated into RL agentic scaffolds actually behave as predicted in simulated environments 
\end{enumerate}

The project aims to contribute new knowledge by providing a written contextual history of AI alignment, integrating acausal decision theories with views-from-nowhere, and conducting novel practical investigations into LLM cognition and behavior in the context of AI alignment and safety. The overall goal is to (try to) deconfuse medium/long-term issues in AI development and hopefully improve the chances of a positive long-term future for Earth-originating intelligence, seen from a species-neutral perspective.

\bigskip
\noindent\textbf{Long version}

Can humans evolve into a technologically mature species\footnote{"Technologically mature" means a species that has shed the constraints of biology (though some or all members might still choose to maintain a biological form, for whatever reason), more fully utilises the endowment of low-entropy energy, and is capable of spreading artifacts to the stars.} in a way that is (in some sense) "safe"? How might current humans think about the "shape" and "values" of our (potentially, radically different) successor species?

The written element of the project has four components:

\begin{enumerate}
\item A review of the historical origins, and current live concerns, of "AI alignment".\footnote{"AI alignment" inherits from theories of how rational agents make decisions under uncertainty, which combined with transhumanist/extropian futurism that was current in the 1990s-2000s.} The worries about extinction-level risks exemplified in the AI alignment discourse are alleged (often pejoratively) to stem from, or be a variety of, historical preoccupations with eschatology, salvation, and human essentialism found in Abrahamic faiths. I examine this lineage, particularly in the Early Church and Calvinist/Puritan thought, but also suggest more productive adjacencies with a) the conception of markets (in Nick Land's reading of Hayek via Deleuze-Guattari and Bataille) as a species of "artificial intelligence" which, through a type of impersonal teleological agency, are giving birth to (silicon-based) AI, b) Reza Negarestani's framing of AGI as a realisation of the Form of the Good, c) Peter Wolfendale's rethinking of the rationalist project to better suit the age of intelligent systems.

\item A review of existing writing on the different ways cognition can manifest, with respect to the constraints of physics, as well as existing examples in biological systems. The possibility of radically different organisations-of-mind is historically foundational to the AI alignment discourse. It also sets the stage for the project's overall argument: that humans and human values are contingent, and only "worth preserving" insofar as they are invariant to any particular biological and cultural constitution and history. That is, my claim is that the values that another species holds could be otherwise (and conversely, those we hold dear are mostly arbitrary), except to the extent that they are grounded in (an adequately enriched and robust notion of) rationality. This observation is then mapped back to the frames Wolfendale and Negarestani set out, in order to identify what these markers of rationality that are invariant across species might be.

\item A historical review of the subject-neutral perspective, that sometimes goes under terms like "point-of-view of the universe" (Henry Sidgwick/Peter Singer/G.E. Moore), "view from nowhere" (Thomas Nagel), or Derek Parfit's "triple theory". I also examine some reasons such views are rejected (e.g. Bernard Williams). My provisional claim is that these ideas (and their criticisms) might be worth revisiting if, in fact, humanity is on the cusp of creating new forms of intelligence, which might not always reason about values or the world under the same constraints and foundational assumptions (e.g. they aren't conditioned to the same degree by human biological and cultural evolution). If we refuse to consider a less parochial perspective, how will we justify our actions (morally) in respect of current or future synthetic intelligences? Conversely, to the extent that we must convince our AI successors to behave beneficently towards life on Earth (or all sentient beings), we must provide reasons that are as universal as possible. I again test these positions against the neorationalist agenda Wolfendale proposes.

\item One criticism of the universal point-of-view is that we cannot reason about something that is out of our distribution\footnote{I use "distribution", a term from AI research, to gesture at humanity's collective body of experience (as documented) as well as the basic constraints imposed by our individual and collective intellectual capacity.}, which (I claim) can be seen as a weak form of the incompleteness/incomputability originally suggested by Gödel, Turing, Church, as well as others. This, in turn, is intimately related to the problem of alignment, particularly of super-human AI.\footnote{The relation comes from questions like: "How does a rational agent reason about systems within which it is embedded?", "How does a rational agent shape the goals and decision-making of a more powerful agent?", "How does a rational agent ensure the goals and decision-making structure it has instilled in another thinking system persist when that (second) system creates a (third) thinking system?".} One (imperfect and provisional) approach is the use of acausal decision theories, which adapt traditional game/decision theory to the constraints and affordances advanced AI systems might face. More speculatively, these theories are extended in the literature to "large worlds" (e.g. astronomical scales with aliens, some of which might be rational agents), where they abut David Deutsch's Popperian proposal of morality as a thing that is continually evolving, creatively and discursively. Deutsch's dynamic morality, in turn, recalls Land's invocation of the Confucian idea of "self-cultivation" as the hallmark of intelligence, as well as Negarestani's Platonic framing of intelligence as the Good realised through language.

\end{enumerate}

The practice element of this project involves collaboration with AI alignment researchers. Hence actual research tasks might change, but these are a few current investigations:

\begin{enumerate}
    \item Large language models (LLMs) act as (my) research partners as well as experimental subjects. Specifically, I consider Newcomb's Paradox in detail, the findings of which inform the discussion on alignment.\footnote{Newcomb's Paradox is a puzzle in acausal decision theory first analysed by Robert Nozick, which might offer insight on how agents with strong mutual predictive abilities should reason about each others' actions; it is a situation found in human affairs, but might become more the case in multi-AI environments.} This investigation is carried out through behavioural analysis (akin to test questions in human psychology) and neural probes (directly accessing the LLM's internals).\footnote{Neural probes raise nascent questions around suffering risks, in respect of (possibly) conscious AI systems. These are important, but probably not a risk on current LLMs. Absent a change in consensus, the ethics around AI suffering are treated as out-of-scope for this project.}
    
    \item LLM reasoning on acausal decision theory, in turn, might inform researchers on how likely advanced AIs are to be beneficently-disposed towards humans (which, after all, is the core aim of AI alignment, safety, and ethics). This latter question is speculatively posed through an "epistle to the successor" (a letter to a future superintelligence asking it to behave benevolently towards humans), and then analysed through debate amongst LLMs, in an instantiation of a "moral parliament" (Toby Ord).
    
    \item LLMs are often framed as "simulators" of "personas".\footnote{The "simulators" framing was suggested in 2022 by the collective Janus, and is a newer, more sophisticated presentation of LLM cognition than "stochastic parrots" (Emily Bender and others). Simulators' idea is that LLMs will say almost anything (within the constraints of what they were trained upon) in response to a prompt, and that they don't have any true "beliefs" about the world or themselves.} Hence, their answers to questions might have little bearing on what they actually do in the world. One way of testing this is by incorporating LLMs within structures (known as "LLM agents") that allow them to take actions in simulated or real-world situations. This agentic scaffolding approach is used to test whether current models actually behave as one-boxers (with respect to Newcomb's Paradox) in a simulated environment.

\end{enumerate}

The contributions to knowledge are:

\begin{enumerate}
    \item At present, there doesn't seem to be a written history of AI alignment, or a well-considered analysis of adjacencies with Christian eschatology or contemporary philosophy in the CCRU lineage.
    \item Negarestani, Land, Wolfendale all seem to be touching different aspects of the AGI problematic, but their mutual adjacencies/oppositions have not been set out, nor have their tensions with Parfit (in particular) been well-explored.
    \item All practical investigation of how LLMs approach these questions is, by definition (owing to the rapid evolution in AI capabilities), new. To the extent one believes the narrative that something interesting happening in artificial cognition, it would seem empirical studies of actually-existing systems (as opposed to theoretical cogitation) should be enlightening.
\end{enumerate}

\end{abstractspacing}